\begin{recipebox}{Coconut Chicken, Potato \& Chickpea Curry}{3--4}{30 min (+ 10 min natural release)}

\begin{ingredients}
  \item 400g chicken breast
  \item 200g tinned tomatoes
  \item 2 tbsp red or green curry paste
  \item 250ml chicken stock
  \item 3 garlic cloves, minced
  \item 1 tsp fresh or ground ginger
  \item 1 tbsp fish sauce or soy sauce
  \item 1 tsp brown sugar or honey
  \item 2 medium potatoes, peeled and diced into 2cm cubes
  \item 1 carrot, grated
  \item 1 can chickpeas, drained and rinsed (400g)
  \item 1 can coconut milk (400ml)
  \item 100g frozen peas
  \item 200g jasmine or basmati rice (dry)
  \item 1 tbsp lime juice (optional)
\end{ingredients}

\begin{method}
  \item Add tinned tomatoes, curry paste, stock, garlic, ginger, fish sauce, and sugar to the Instant Pot insert. Stir to combine --- keep the mix thin and well stirred so nothing thick sits on the bottom.
  \item Place the whole chicken breast into the liquid, then scatter the diced potato around it. Do \textbf{not} stir again after this --- leave the base layer undisturbed to avoid a burn warning.
  \item Seal the lid, set valve to Sealing, and cook on Manual/Pressure Cook on HIGH for \textbf{10 minutes}.
  \item Let pressure release naturally for 8--10 minutes, then carefully quick release any remaining pressure.
  \item Remove the chicken breast and shred it with two forks. Stir in the chickpeas and return the shredded chicken to the pot.
  \item Switch to Saut\'{e} mode. Stir in the coconut milk, grated carrot, and frozen peas. Heat through for 3--5 minutes, stirring occasionally, until everything is hot and the peas are cooked.
  \item Turn off Saut\'{e}. Add lime juice if using and taste for seasoning --- more fish sauce for salt, more sugar to balance.
  \item Cook the rice according to packet instructions (\textasciitilde15 min, or use the Instant Pot's rice function separately) and serve with the curry.
\end{method}

\begin{notes}
\textbf{Why the coconut milk goes in last:} Coconut milk and the thicker tomato/paste mix sitting on the bottom can trigger a burn warning under pressure. Keeping the base liquid thin (stock-heavy) and stirring the coconut milk in afterwards on Saut\'{e} avoids this.

\textbf{Grated carrot:} Goes in with the coconut milk at the end rather than under pressure --- it cooks through in a couple of minutes on Saut\'{e} and keeps a bit of texture and sweetness.

\textbf{Chicken breast:} 10 minutes at high pressure plus the natural rest is enough for a single 400g breast --- don't add extra time or it will dry out. If cooking from frozen, add 3--4 minutes.
\end{notes}

\end{recipebox}
